
``Tempo seco e altas temperaturas são fatores que favorecem o aumento de incêndios ambientais'' \cite{Silva21}. Em tempos como os atuais, faz-se ainda mais necessário encontrar ferramentas que auxiliem a frenagem destes eventos. Recentemente, o Instituto Nacional de Metereologia (INMET) emitiu diversos alertas sobre as ondas de calor que viriam a atingir algumas regiões do país. Como consequência o segundo semestre do ano de 2023, em especial os meses de agosto à novembro, registraram um recorde histórico na média de temperaturas, conforme exposto na reportagem publicada pelo próprio instituto, ``Temperatura média atinge recorde no Brasil pelo quarto mês seguido'' \cite{inmet23:ondadecalor}.

Os prejuízos causados por incêndios e queimadas, principalmente quando se trata da Amazônia, que é considerada por muitos como o pulmão do mundo, são diversos. Silva (\citeyear{Silva21}) argumenta como estes eventos introduzem um processo de degradação do solo, perda de matéria-prima e poluição do ar, desencadeando impactos devastadores que reverberam além das fronteiras da floresta, afetando ecossistemas globais e sendo gerador de prejuízos econômicos com uma magnitude preocupante.

Estes danos, apesar de já se mostrarem severos, podem ir muito além. Muitos estudos evidenciam que a nossa participação na emissão de gases, como o dióxido de carbono, alavancam o fenômeno de aquecimento global. Existem, claro, opiniões sobre o tema que se contrapõem e devem ser consideradas. Segundo a reportagem ``Aquecimento global: quem é culpado?" \cite{Aguiar13}, o efeito estufa é uma realidade inegável, mas a dúvida paira sobre a real contribuição do homem neste âmbito. Caso o aquecimento global antropogênico seja factual, as incidências de queimadas corroboram para o cenário atual, o que torna ainda mais importante o estudo e a prevenção das mesmas.

A utilização de agentes inteligentes (como máquinas de aprendizado) para este escopo de pesquisa representa uma abordagem promissora no campo da gestão ambiental e prevenção de desastres. Os modelos de aprendizado de máquina são sistemas computacionais inspirados na racionalidade humana, capazes de aprender padrões complexos a partir de dados. Ao serem aplicadas na previsão de incêndios florestais, esses agentes podem ter a capacidade de assimilar grandes volumes de informações (por exemplo, no espectro de Redes Neurais Artificiais), considerando variáveis ambientais, climáticas e históricas. De acordo com o modelo utilizado pode-se ter flexibilidade na modelagem da sua arquitetura, o que permite a adaptação a padrões não-lineares presentes nos dados, objeto fundamental para modelar a dinâmica complexa dos incêndios florestais \cite{Ferneda06}.

Treinar um modelo com conjuntos de dados históricos que contêm informações detalhadas sobre a topologia e vegetação de uma região é um desafio desde o tratamento dos dados até a avaliação dos mesmos. A realização deste projeto busca desenvolver uma inteligência apta a reconhecer similaridades sutis e correlações que podem ser cruciais na identificação das áreas propensas a incêndios. O uso desse tipo de tecnologia pode ser útil às unidades competentes, contribuindo com o enfrentamento da degradação ambiental por meio de um mapa de risco capaz de orientar sobre as áreas que possuem mais suscetibilidade a combustão. Assim, a proposta é ser capaz de obter um sistema eficiente e de baixo custo computacional para a previsão de incêndios em um ambiente com uma grande quantidade de dados disponíveis.

\section{Justificativa}

O monitoramento de incêndios florestais é uma tarefa que transita entre os espectros de previsão sobre o risco de fogo e a detecção em tempo real. Prever é frequentemente considerado mais interessante ou valioso do que simplesmente detectá-los após o início de uma incidência, por várias razões. Primeiro, permite que medidas preventivas sejam tomadas antes mesmo do início do fogo. Isso pode incluir a alocação mais eficiente de recursos de combate a incêndios e decisões de gestão florestal a longo prazo, como o planejamento de zonas de proteção em torno de áreas sensíveis.

Além disso, prever incêndios pode simplesmente ajudar a reduzir as perdas econômicas associadas a danos à propriedade, perda de colheitas e interrupção de atividades comerciais. Isso é especialmente importante em áreas onde a agricultura, o turismo e outras atividades econômicas dependem da saúde dos ecossistemas naturais. Do outro lado da moeda, um dos principais benefícios de sistemas de detecção é a capacidade de identificar incêndios em estágios iniciais. Isso permite uma resposta rápida e eficaz para conter o fogo antes que ele se espalhe e cause danos significativos. Em suma, pode-se dizer que tanto a detecção quanto a previsão de incêndios são importantes; a capacidade de detectá-los fornece uma resposta rápida e eficaz, enquanto a de prevê-los oferece uma gama mais ampla de benefícios, desde a prevenção completa de danos até a redução de perdas econômicas.

A urgência do tema é corroborada por publicações recentes na literatura internacional. Aragão \emph{et al.}~\cite{aragao2018amazon} documentaram, no periódico \emph{Nature Communications}, que desde o início do século~XXI o fogo associado à seca passou a \emph{neutralizar} os ganhos obtidos com a redução do desmatamento na Amazônia. Mais recentemente, Silva-Junior \emph{et al.}~\cite{silvajunior2025amazon}, em \emph{Biogeosciences}, mostraram que 2024 marcou a pior degradação por fogo em mais de duas décadas: \SI{3,3}{\mega\hectare} queimados (acréscimo de \(\approx 400\%\) em relação aos dois anos anteriores), com emissão estimada de \SI{791}{\mega\tonne} de CO$_{2}$. Esses números reforçam a urgência operacional do desenvolvimento de ferramentas modernas de previsão e vigilância, capazes de produzir mapas de risco em tempo quase-real e auditáveis por gestores ambientais.

Dada a magnitude e complexidade das queimadas na região da Amazônia, abordagens eficientes para prevenção, monitoramento e resposta rápida são cruciais. A utilização de agentes inteligentes para mapear focos de incêndio não é uma novidade, contudo, dadas as especificidades de cada região atacada, diferentes variáveis podem se destacar nos algoritmos para a previsão dos focos. Além disso, deve-se considerar o contexto global das mudanças climáticas e o impacto direto que possuem as queimadas sobre a vegetação e biodiversidade amazônica nessas alterações. Isto posto, este projeto deve instituir uma base de conhecimento sólida, integrando conceitos de sensoriamento remoto e aprendizado de máquinas na intenção de entender e mitigar os efeitos adversos dos incêndios florestais.


\section{Objetivos}

\subsection{Objetivo Geral}

Desenvolver um mapa das regiões suscetíveis a incêndios florestais na área da Amazônia Legal, utilizando técnicas de aprendizado de máquina. O projeto tem como propósito oferecer contribuições relevantes para a previsão e o monitoramento de queimadas, além de classificar as regiões com maior propensão à ocorrência desses eventos.

\subsection{Objetivos Específicos}

\begin{itemize}
    \item[1.] Realizar uma revisão bibliográfica sobre as técnicas de aprendizado de máquina e classificação em evidência na literatura, incluindo trabalhos publicados entre 2024 e 2025 que utilizam \emph{ensembles} e métodos de interpretabilidade aplicados a risco de incêndio florestal.
    \item[2.] Mapear as regiões da Amazônia Legal que sofreram com queimadas durante a última década (2014--2023) e definir suas principais causas, com base em registros do Programa Queimadas/INPE.
    \item[3.] Obter e enriquecer os dados climáticos, topológicos e de incêndios históricos, incluindo enriquecimento com NASA POWER (umidade relativa RH2M) e NASA FIRMS (FRP em tempo quase-real), além da construção de \textbf{24~\emph{features} físico-climáticas avançadas} (Tier~1: SPI, KBDI \emph{proxy}, VPD \emph{proxy}, anomalias de precipitação, histórico estendido de fogo) embasadas na literatura recente~\cite{seager2015climatology,forests2024cerradoamazon,npjhazards2025fire}.
    \item[4.] Desenvolver um algoritmo que utilize diferentes modelos de aprendizado de máquina para a previsão de regiões mais suscetíveis a queimadas, comparando oito abordagens (Regressão Logística, SGD, Random Forest com SMOTE, RF balanceado otimizado por \textbf{Optuna}~\cite{akiba2019optuna}, XGBoost~\cite{chen2016xgboost}, LightGBM~\cite{ke2017lightgbm}, CatBoost~\cite{prokhorenkova2018catboost}, e \emph{ensembles} \emph{Voting Soft} e \emph{Stacking}~\cite{wolpert1992stacked}), com calibração isotônica de probabilidades~\cite{zadrozny2002transforming} e ajuste multi-classe de limiares de decisão.
    \item[5.] Avaliar a eficácia dos modelos propostos por meio de métricas de desempenho por classe (precisão, revocação, \(\text{F}_1\), \(\text{F}_1\)-macro), validação cruzada estratificada \(k\)-\emph{fold} e \textbf{validação temporal estrita} (\emph{rolling-origin} por ano)~\cite{bergmeir2012cv} para quantificar o viés do \emph{split} aleatório, complementadas por análise de interpretabilidade via \textbf{SHAP}~\cite{lundberg2017unified,lundberg2020local} e \textbf{\emph{Permutation Importance}}~\cite{breiman2001random,fisher2019all}.
    \item[6.] Desenvolver uma \textbf{aplicação web interativa explicável} (\emph{Flask} + \emph{Folium}~\cite{folium,leafletjs}) que, ao receber um clique do usuário sobre qualquer ponto da Amazônia Legal, classifique o risco do ponto consultado e exiba a explicação local das \emph{features} responsáveis pela decisão, viabilizando uso operacional por gestores ambientais.
\end{itemize}

\section{Estrutura da Monografia}

Este trabalho está organizado em seis capítulos que abordam diferentes etapas da pesquisa, proporcionando uma compreensão abrangente do desenvolvimento do sistema de previsão de queimadas na Amazônia.

\textbf{2. Estado da Arte:} são explorados trabalhos correlatos e avanços significativos na área de previsão de queimadas, com foco em estudos que propõem o mapeamento de regiões suscetíveis a fogo. A revisão aborda a extensão do problema considerando seus impactos ambientais, sociais e econômicos. Adicionalmente, são examinadas abordagens tradicionais de aprendizado de máquina e o início dos sistemas de prevenção, complementadas por uma seção de \textbf{literatura recente 2024--2025} (\emph{ensembles} explicáveis, índice KBDI no Cerrado-Amazônia, validação temporal estrita).

\textbf{3. Fundamentação Teórica:} estabelece as bases conceituais necessárias para a compreensão do problema das queimadas na Amazônia. Inicia-se com a contextualização do bioma amazônico e a problemática dos incêndios florestais, seguida pela discussão do sensoriamento remoto e das tecnologias espaciais. O capítulo prossegue com os fundamentos da inteligência artificial e do aprendizado de máquina, contemplando \textbf{modelos clássicos} (Regressão Logística, Random Forest), \textbf{modelos modernos de \emph{boosting}} (XGBoost, LightGBM, CatBoost), \textbf{\emph{ensembles}} (Voting e Stacking), \textbf{otimização de hiperparâmetros} (Optuna), \textbf{interpretabilidade} (SHAP e Permutation Importance), \textbf{tratamento de desbalanceamento}, \textbf{calibração de probabilidades}, \textbf{índices físico-climáticos de fogo} (KBDI, SPI, VPD, De Martonne) e \textbf{validação temporal} de modelos.

\textbf{4. Metodologia:} descreve detalhadamente os passos e procedimentos adotados para desenvolver o sistema de previsão. São delineados os métodos de coleta e enriquecimento dos dados (Programa Queimadas + NASA POWER + NASA FIRMS), o pré-processamento, a engenharia das 24~\emph{features} Tier~1, o protocolo de treinamento dos modelos, a avaliação (métricas, interpretabilidade, \emph{threshold tuning}, validação temporal), e o desenvolvimento da aplicação web interativa explicável.

\textbf{5. Resultados:} apresenta os resultados experimentais: comparativo entre as oito abordagens treinadas, desempenho detalhado do modelo final, análise de interpretabilidade por SHAP e \emph{Permutation Importance}, resultados da validação temporal estrita, efeito do \emph{threshold tuning}, evolução incremental do desempenho ao longo das decisões metodológicas, validação operacional via consultas comparativas no aplicativo web e \textbf{estudo de caso externo} (Pium--TO / CENSIPAM, fusão INMET WIS2/ZIP, extensões M1--M4).

\textbf{6. Conclusão:} sintetiza os principais achados, registra as limitações reconhecidas e enumera os trabalhos futuros prioritários para evoluir o sistema rumo à integração de novas fontes (NDVI/EVI MODIS, MapBiomas Fogo) e de causalidade estrita nas \emph{features} de janela móvel.

\chapter{Estado da Arte}
